\documentclass[12pt]{article}
\usepackage{amsmath, amssymb, amsthm}
\usepackage{geometry}
\usepackage{hyperref}
\usepackage{bm}
\usepackage{booktabs}
\usepackage{setspace}
\usepackage{xcolor}
\usepackage{mdframed}
\usepackage{enumitem}
\onehalfspacing
\geometry{margin=1.2in}

\newtheorem{lemma}{Lemma}
\newtheorem{proposition}{Proposition}
\newtheorem{remark}{Remark}

\newmdenv[backgroundcolor=blue!5, linecolor=blue!30, roundcorner=4pt]{insight}
\newmdenv[backgroundcolor=orange!5, linecolor=orange!30, roundcorner=4pt]{codebox}
\newmdenv[backgroundcolor=green!5, linecolor=green!30, roundcorner=4pt]{keyresult}

\title{How Welfare Loss Is Computed \\ \large A Teaching Note: Derivation, a Worked Toy Example, and How It Scales to the Real Chile Model}
\author{Working Notes}
\date{\today}

\begin{document}
\maketitle
\tableofcontents
\newpage

%==============================================================
\section{The Object We're Computing}
%==============================================================

Second-order welfare loss relative to the efficient (flex-price) allocation, from Rubbo's
Proposition 3 as extended in this project (full proof in
\texttt{rubbo\_proofs\_and\_extension.tex}, \S``Welfare Function''):
\begin{equation}
    \mathbb{W} = \frac{1}{2}\,\mathbb{E}\left[\underbrace{(\gamma+\varphi)\,y_t^2}_{\text{output gap}}
    \;+\; \underbrace{\sum_i \kappa_i\, \pi_{it}^2}_{\text{within-sector}}
    \;+\; \underbrace{\Phi_C(\bm\mu_t,\bm\mu_t) + \sum_s\lambda_s\Phi_s(\bm\mu_t,\bm\mu_t)}_{\text{cross-sector}}\right]
    \label{eq:welfare}
\end{equation}
where $\kappa_i \equiv \lambda_{D,i}\,\varepsilon_i\,\frac{1-\hat\delta_i}{\hat\delta_i}$ is sector
$i$'s \emph{welfare weight} — large when the sector is big in the network ($\lambda_{D,i}$) and
sticky ($\hat\delta_i$ small, so any inflation it does have translates into a lot of costly
within-sector price dispersion). $\bm\mu_t$ is the vector of log markup gaps.

\begin{insight}
\textbf{Why three terms, not one.} A representative-agent closed economy only has the output-gap
term. Once you add (i) multiple sectors with sticky prices, each sector generates its own
Calvo price-dispersion cost — Term 2; and (ii) non-unit elasticities of substitution across
goods, relative-price misallocation across sectors is itself costly even holding each sector's own
dispersion fixed — Term 3. Term 3 is usually the smallest (in the Chile calibration it adds only
3--8\% on top of Terms 1+2 — see \S\ref{sec:real}) but it's not structurally negligible; it's small
here because this model's production nests are Cobb-Douglas (elasticity exactly 1), which is a
property of the functional form, not a generic result.
\end{insight}

Because $\mathbb{W}$ is a sum of expectations of \emph{squares} of variables that are themselves
linear in exogenous shocks (in a linearized model), $\mathbb{W}$ reduces entirely to a weighted sum
of \textbf{variances and covariances}. This is the fact that makes every decomposition below exact,
not approximate:
\begin{equation}
    y_t = \sum_k L_{y,k}\,\epsilon_{k,t}, \qquad \pi_{it} = \sum_k L_{i,k}\,\epsilon_{k,t}
    \label{eq:linear}
\end{equation}
for structural shocks $\epsilon_{k,t}$ with $\mathrm{Cov}(\epsilon_k,\epsilon_{k'})=0$ for
$k\neq k'$ (the model's shocks are mutually independent by construction). Then
$\mathrm{Var}(\pi_i) = \sum_k L_{i,k}^2\,\sigma_k^2$ — additive across shocks, with no
cross-terms, \emph{precisely because independence kills the covariance terms}. This is the whole
trick behind "shock $k$ contributes $X\%$ of the welfare loss."

%==============================================================
\section{Full Derivation: From Utility to Welfare Loss}
\label{sec:derivation}
%==============================================================

Eq.~\eqref{eq:welfare} isn't an assumption — it's what a second-order Taylor expansion of household
utility around the efficient allocation produces, once you account for Calvo price stickiness and
CES/Cobb-Douglas aggregation. This section derives each of the three terms in turn, at a level of
detail one notch below \texttt{rubbo\_proofs\_and\_extension.tex}'s ``Proposition 3'' proof outline
(\S6 there) — the goal here is to make every algebraic step explicit rather than just stating the
result.

\subsection{Term 1: the output-gap term, from household utility directly}

Household period utility is $U(C_t,L_t) = u(C_t) - v(L_t)$ with $u(C)=\frac{C^{1-\gamma}}{1-\gamma}$,
$v(L)=\frac{L^{1+\varphi}}{1+\varphi}$. Let $(C_t^*, L_t^*)$ denote the efficient (flex-price,
correctly-subsidized) allocation, and write $\tilde c_t = \log(C_t/C_t^*)$,
$\tilde l_t = \log(L_t/L_t^*)$ for log-deviations from it.

\textbf{Expand $u(C_t)$ to second order.} Writing $C_t = C_t^* e^{\tilde c_t}$ and Taylor-expanding
the exponential to second order in $\tilde c_t$:
\begin{equation}
    u(C_t) - u(C_t^*) = \frac{(C_t^*)^{1-\gamma}}{1-\gamma}\Big[(1-\gamma)\tilde c_t + \tfrac{1}{2}(1-\gamma)^2\tilde c_t^2\Big] + O(\tilde c_t^3)
    = U_C^*C_t^*\Big[\tilde c_t + \tfrac{1}{2}(1-\gamma)\tilde c_t^2\Big]
\end{equation}
using $U_C^* \equiv u'(C_t^*) = (C_t^*)^{-\gamma}$, so $U_C^*C_t^* = (C_t^*)^{1-\gamma}$. The
symmetric calculation for $-v(L_t)$ (with $V_L^* \equiv v'(L_t^*) = (L_t^*)^\varphi$) gives
\begin{equation}
    -v(L_t) + v(L_t^*) = -V_L^*L_t^*\Big[\tilde l_t + \tfrac{1}{2}(1+\varphi)\tilde l_t^2\Big]
\end{equation}

\textbf{Adding them:}
\begin{equation}
    \Delta U_t \equiv U_t - U_t^* = U_C^*C_t^*\tilde c_t - V_L^*L_t^*\tilde l_t
    \;+\; \tfrac{1}{2}U_C^*C_t^*(1-\gamma)\tilde c_t^2 - \tfrac{1}{2}V_L^*L_t^*(1+\varphi)\tilde l_t^2
    \label{eq:deltaU}
\end{equation}

\textbf{Killing the linear terms.} Two facts collapse the first (linear) bracket to exactly zero:
\begin{enumerate}[label=(\roman*)]
    \item \textbf{Goods-market clearing to first order:} $\tilde c_t = \tilde y_t$.
    \item \textbf{Lemma 2 (output gap = employment gap, \texttt{rubbo\_proofs\_and\_extension.tex} \S4):} $\tilde y_t = \tilde l_t$, because sticky prices only distort the \emph{aggregate} labor margin — the allocation of a given aggregate $L_t$ across sectors stays efficient since relative wages are flexible.
\end{enumerate}
So $\tilde c_t = \tilde l_t = \tilde y_t$ \emph{exactly} at this order. Combined with the efficient
labor-supply condition $U_C^*C_t^* = V_L^*L_t^*$ (the standard welfare-theorem normalization: at the
efficient point, the household's MRS equals the marginal product, which after the usual optimal
employment subsidy correcting the monopoly markup makes the two curvature-weighted levels equal),
the linear term is
\begin{equation}
    U_C^*C_t^*\tilde c_t - V_L^*L_t^*\tilde l_t = U_C^*C_t^*(\tilde y_t - \tilde y_t) = 0
\end{equation}
\textbf{exactly} — not approximately, because $\tilde c_t$ and $\tilde l_t$ are the \emph{same}
variable here, not merely correlated.

\textbf{What's left} is purely the quadratic terms in Eq.~\eqref{eq:deltaU}, all in $\tilde y_t$ (via $\tilde c_t=\tilde l_t=\tilde y_t$) and all sharing the common factor $U_C^*C_t^*$:
\begin{equation}
    \Delta U_t = \tfrac{1}{2}U_C^*C_t^*\Big[(1-\gamma) - (1+\varphi)\Big]\tilde y_t^2
    = -\tfrac{1}{2}(\gamma+\varphi)\,U_C^*C_t^*\,\tilde y_t^2
\end{equation}
Dividing through by $U_C^*C_t^*$ (the standard normalization to express welfare loss in
consumption-equivalent units) and flipping sign to define $\mathbb{W}\equiv -\Delta U_t/U_C^*C_t^*
\geq 0$ gives exactly Term 1 of Eq.~\eqref{eq:welfare}: $\tfrac12(\gamma+\varphi)\tilde y_t^2$.

\begin{insight}
The output-gap term needed \emph{no} information about the production network, Calvo parameters, or
number of sectors — it comes purely from curvature of preferences ($\gamma,\varphi$) and the two
identities $\tilde c_t=\tilde y_t=\tilde l_t$. All network/price-stickiness structure enters only
through Terms 2 and 3, and indirectly, through how $\tilde y_t$ itself is determined in equilibrium
(the Phillips curve system).
\end{insight}

\subsection{Term 2: within-sector price dispersion, from Calvo pricing}

Within sector $i$, firms that reset at different past dates charge different prices $P_{if,t}$, even
though they face identical current marginal cost. This dispersion is directly costly: producing a
given amount of the sector-$i$ CES aggregate $Y_{it}=\big(\int_0^1 Y_{if,t}^{(\varepsilon_i-1)/\varepsilon_i}df\big)^{\varepsilon_i/(\varepsilon_i-1)}$
from dispersed-price varieties wastes labor/inputs relative to charging one uniform price, because
demand $Y_{if}=(P_{if}/P_i)^{-\varepsilon_i}Y_i$ shifts disproportionately toward the (arbitrarily)
cheaper varieties without them being any more efficient to produce.

\textbf{The dispersion wedge.} Define $\Delta_{it} \equiv \int_0^1 (P_{if,t}/P_{it})^{-\varepsilon_i}\,df \geq 1$
(equality iff all firms charge the same price). A second-order expansion of $\Delta_{it}$ in the
cross-sectional log-price distribution gives the standard result (Yun 1996; Woodford 2003, ch.~6):
\begin{equation}
    \log \Delta_{it} \approx \tfrac{1}{2}\varepsilon_i\,\mathrm{Var}_f[\log P_{if,t}]
    \label{eq:dispersion_wedge}
\end{equation}
— a pure Jensen's-inequality cost of the CES aggregator, increasing in both the elasticity
$\varepsilon_i$ (more substitutable varieties $\Rightarrow$ demand reallocates more aggressively
toward mispriced ones) and the cross-sectional price variance.

\textbf{Mapping dispersion to inflation.} Under Calvo, $\mathrm{Var}_f[\log P_{if,t}]$ is not a free
parameter — it's pinned down by the sector's own inflation history, because dispersion only arises
from \emph{staggering}: firms that reset at $t$ all charge the identical $P_{it}^*$, so all
cross-sectional variance comes from the gap between resetters and the (persisting) distribution of
past prices. Log-linearizing the recursion $\mathrm{Var}_f[\log P_{if,t}] = (1-\delta_i)\big(\mathrm{Var}_f[\log P_{if,t-1}] + (\text{new dispersion from this period's inflation})\big)$
around the zero-inflation steady state and solving forward (Yun 1996; the algebra is a standard but
lengthy application of the Calvo recursion, omitted here) yields
\begin{equation}
    \mathrm{Var}_f[\log P_{if,t}] \approx \frac{1-\hat\delta_i}{\hat\delta_i}\,\pi_{it}^2
    \label{eq:calvo_var}
\end{equation}
with $\hat\delta_i$ the same effective Calvo parameter defined in
\texttt{rubbo\_proofs\_and\_extension.tex} Eq.~(19) (``Calvo Optimal Reset Price''). Intuitively: a stickier sector ($\hat\delta_i$ small) needs a \emph{larger} cross-sectional
variance to generate the same observed aggregate inflation $\pi_{it}$, because fewer firms are
moving to produce that inflation number, so the ones that do move must move further apart from the
frozen majority.

\textbf{From dispersion to a TFP-equivalent utility loss.} $\log\Delta_{it}$ acts exactly like a
negative TFP shock to sector $i$: producing sector output $Y_{it}$ requires
$\Delta_{it}\times$(what would be needed absent dispersion) units of underlying labor/inputs. Its
contribution to aggregate consumption-equivalent loss is this TFP loss weighted by the sector's
\textbf{Domar weight} $\lambda_{D,i}$ (its total-sales share, counting indirect sales through the
network) — the same weighting Lemma 1 uses to build \emph{natural} output from sectoral TFP:
\begin{equation}
    \text{Term 2} = \sum_i \lambda_{D,i}\,\varepsilon_i \cdot \tfrac{1}{2}\mathrm{Var}_f[\log P_{if,t}]
    \overset{\eqref{eq:calvo_var}}{=} \tfrac12\sum_i \underbrace{\lambda_{D,i}\,\varepsilon_i\,\frac{1-\hat\delta_i}{\hat\delta_i}}_{\displaystyle \kappa_i}\,\pi_{it}^2
\end{equation}
exactly Term 2 of Eq.~\eqref{eq:welfare}.

\subsection{Term 3: cross-sector misallocation, from non-unit substitution elasticities}

Terms 1--2 together would be the \emph{entire} welfare loss if every aggregator in the model
(household consumption bundle, each firm's input bundle) were Cobb-Douglas — because Cobb-Douglas
has unit elasticity of substitution, and a firm/household with unit elasticity is, to second order,
\emph{indifferent} to small relative-price distortions across its inputs (it substitutes away from
the relatively-expensive input exactly enough to keep expenditure shares fixed, which is precisely
what makes the resulting misallocation cost vanish at second order). Term 3 exists only because this
model's \emph{outer} nest — $C_t=\mathrm{CES}(C_t^H,C_t^F;\eta)$, domestic vs.\ imported consumption
— has $\eta=1.5\neq1$.

\textbf{General second-order misallocation cost (Baqaee--Farhi).} For any CES-type aggregator
$G(X_1,\ldots,X_n)$ with cost shares $w_k$ and pairwise Allen--Uzawa elasticities $\sigma_{kh}$, a
vector of small \emph{wedges} $x_k$ (log deviations of $X_k$ from its cost-minimizing/efficient
level) generates a second-order TFP-equivalent loss
\begin{equation}
    \Phi_G(\bm{x},\bm{x}) = \tfrac{1}{2}\sum_{k,h} w_k w_h\,\sigma_{kh}\,(x_k-x_h)^2
    \label{eq:phi_general}
\end{equation}
— a Harberger-triangle-style formula: quadratic in the wedge \emph{differences} (only relative
distortions matter, not the level), scaled by how substitutable the pair is. This is the general
form; deriving it from a second-order expansion of $G$ itself is standard in the
production-network literature (Baqaee \& Farhi 2019, 2020) and not reproduced here.

\textbf{The wedges here are markup gaps.} In this model, the object playing the role of $x_k$ is the
log markup gap $\mu_{it} = \log(P_{it}/MC_{it}) - \log\frac{\varepsilon}{\varepsilon-1}$ — the
deviation of sector $i$'s markup from its flexible-price target — applied once at the household's
consumption nest ($\Phi_C$, weights $\bm\beta^H$) and once per firm's production nest ($\Phi_s$,
weights $(\alpha_s,\Omega^H_{s\cdot},\Omega^F_s)$), summed with Domar weights on the latter:
\begin{equation}
    \text{Term 3} = \Phi_C(\bm\mu_t,\bm\mu_t) + \sum_s \lambda_{D,s}\,\Phi_s(\bm\mu_t,\bm\mu_t)
\end{equation}
exactly Term 3 of Eq.~\eqref{eq:welfare}. The closed-form collapse to a \emph{weighted variance} of
$\bm\mu_t$ used in \S\ref{sec:real} (and in \texttt{code/cross\_sector\_welfare.py}) follows from
Eq.~\eqref{eq:phi_general} whenever $\sigma_{kh}\equiv\sigma$ is constant across all pairs — true
here since $\sigma_{kh}=1$ (Cobb-Douglas) for every domestic-domestic and domestic-labor-import pair,
and only the domestic-vs.-import-composite pairs get $\sigma_{kh}=\eta$.

\begin{keyresult}
\textbf{Why the three terms have the shape they do.} Term 1 comes purely from preference curvature
and holds with \emph{no} reference to price stickiness at all (it would exist even with flexible
prices, if output happened to deviate from efficient for some other reason). Terms 2 and 3 both come
from the \emph{same} underlying source — CES aggregation of heterogeneous prices — but at two
different levels of the model: Term 2 is dispersion \emph{within} a sector's own variety continuum
(entirely a Calvo/stickiness artifact, vanishes if $\delta_i=1$), while Term 3 is misallocation
\emph{across} sectors/goods in a bundle (a property of the aggregator's curvature $\sigma_{kh}$,
vanishes only if every relevant elasticity is exactly 1, regardless of stickiness). This is also
why Term 3 needed a genuinely new elasticity object ($\eta$) that Terms 1--2 never used.
\end{keyresult}

%==============================================================
\section{Worked Toy Example}
%==============================================================

To make this concrete before touching the real 3-sector Chile model, consider a minimal toy economy
with \textbf{two} sectors: Sector 1 (upstream, directly FX-exposed) and Sector 2 (downstream,
buys from Sector 1, resets prices slowly). Three independent shocks hit each period: $\chi_1,\chi_2$
(sector TFP) and $\Delta e$ (FX depreciation). Full code: \texttt{toy\_welfare.py} (scratchpad).

\subsection{Primitives}
\begin{center}
\begin{tabular}{lcc}
\toprule
 & Sector 1 (upstream) & Sector 2 (downstream) \\
\midrule
Domar weight $\lambda_i$ & 0.30 & 0.70 \\
Calvo reset prob. $\hat\delta_i$ & 0.40 & 0.15 \\
Welfare weight $\kappa_i = \lambda_i\varepsilon\frac{1-\hat\delta_i}{\hat\delta_i}$ & 2.70 & 23.80 \\
\bottomrule
\end{tabular}
\end{center}
($\varepsilon=6$, $\gamma+\varphi=1$ for both sectors — Sector 2's $\kappa$ is nearly 9$\times$
Sector 1's despite being only $2.3\times$ larger in Domar weight, because it is also far stickier.)

\subsection{Impact loadings (a stand-in for the full Phillips-curve solution)}
Instead of solving a full DSGE, assign each variable a linear loading on each shock directly —
this is exactly what the real model's $\mathcal{B}, \mathcal{V}, \Gamma$ vectors do (network-amplified
loadings), just without deriving them from IO data:
\begin{center}
\begin{tabular}{lccc}
\toprule
 & $\chi_1$ & $\chi_2$ & $\Delta e$ \\
\midrule
$\pi_1$ & 1.00 & 0.00 & 0.25 \\
$\pi_2$ & 0.40 & 1.00 & 0.08 \\
$y$     & 0.50 & 0.50 & $-0.20$ \\
$\mu_1$ & 0.50 & 0.00 & 0.12 \\
$\mu_2$ & 0.20 & 0.50 & 0.04 \\
\bottomrule
\end{tabular}
\end{center}
Sector 2's loading of $0.40$ on $\chi_1$ (a shock to Sector \emph{1}) is the network channel: Sector 2
inherits cost-push from its upstream supplier even though $\chi_1$ never touches it directly — the
same logic as $\mathcal{V}=\Delta(I-\Omega\Delta)^{-1}(\cdots)$ in the real model. Shock sizes:
$\sigma_{\chi_1}=\sigma_{\chi_2}=0.010$, $\sigma_{\Delta e}=0.020$ (FX more volatile, as calibrated
in the real model).

\subsection{Results}

\textbf{Term decomposition} (Eq.~\ref{eq:welfare}, evaluated numerically):
\begin{center}
\begin{tabular}{lrr}
\toprule
Term & Value ($\times10^{-3}$) & Share \\
\midrule
Output-gap & 0.033 & 2.0\% \\
Within-sector & 1.580 & 97.7\% \\
\quad Cross-sector & 0.0038 & 0.2\% \\
\midrule
\textbf{Total} $\mathbb{W}$ & \textbf{1.616} & 100\% \\
\bottomrule
\end{tabular}
\end{center}
Within-sector dispersion dominates — same qualitative pattern as the real Chile result
(\S\ref{sec:real}), and for the same structural reason: Cobb-Douglas cross-sector nests make the
cross-sector term small \emph{by construction} (elasticity exactly 1), not because misallocation is
empirically unimportant.

\textbf{Sector decomposition} of the within-sector term:
\begin{center}
\begin{tabular}{lrr}
\toprule
 & Contribution ($\times10^{-3}$) & Share of within-sector \\
\midrule
Sector 1 (upstream) & 0.169 & 10.7\% \\
Sector 2 (downstream/sticky) & 1.411 & 89.3\% \\
\bottomrule
\end{tabular}
\end{center}
Despite Sector 1 having \emph{larger} exposure to the FX shock directly, Sector 2 absorbs almost
90\% of the within-sector loss — because $\kappa_2 \gg \kappa_1$ dominates $\mathrm{Var}(\pi_1) >
\mathrm{Var}(\pi_2)$. This is the toy version of the real model's ``Why Services?'' result: size
$\times$ stickiness beats direct exposure.

\textbf{Shock decomposition} of total $\mathbb{W}$ (exact, by additivity of Eq.~\ref{eq:linear}):
\begin{center}
\begin{tabular}{lrr}
\toprule
Shock & Contribution ($\times10^{-3}$) & Share \\
\midrule
$\chi_1$ (Sector-1 TFP) & 0.339 & 21.0\% \\
$\chi_2$ (Sector-2 TFP) & 1.206 & 74.6\% \\
$\Delta e$ (FX depreciation) & 0.073 & 4.5\% \\
\bottomrule
\end{tabular}
\end{center}
Even though the FX shock is \emph{twice} as volatile as either TFP shock, it contributes the least
to welfare loss — because neither sector's $\kappa_i$ nor its $\Gamma_i$-type loading on $\Delta e$
is large enough to compound with that extra volatility. This is exactly the kind of result that
looks surprising until you decompose it, and exactly why the real project's risk-premium-volatility
sweep (\S\ref{sec:regime}) was necessary rather than assumed.

\textbf{Between/within split of the within-sector term.} Term 2 itself further decomposes, exactly
like weighted ANOVA, into an \emph{aggregate}-inflation piece and a \emph{dispersion}-around-that-
aggregate piece:
\begin{equation}
    \sum_i \kappa_i \pi_{it}^2 \;=\; \kappa_{\mathrm{tot}}\,\bar\pi_t^2 \;+\; \kappa_{\mathrm{tot}}\sum_i w_i(\pi_{it}-\bar\pi_t)^2,
    \qquad w_i \equiv \kappa_i/\kappa_{\mathrm{tot}},\ \ \bar\pi_t\equiv\sum_i w_i\pi_{it}
    \label{eq:between_within}
\end{equation}
an identity ($E[X^2]=(E[X])^2+\mathrm{Var}(X)$ applied with $\kappa$-weights), true state-by-state,
not just on average. In the toy example:
\begin{center}
\begin{tabular}{lrr}
\toprule
 & Value ($\times10^{-3}$, raw, no $\tfrac12$) & Share of Term 2 \\
\midrule
Between ($\kappa_{\mathrm{tot}}\bar\pi^2$) & 2.801 & 88.7\% \\
Within (dispersion around $\bar\pi$) & 0.358 & 11.3\% \\
\bottomrule
\end{tabular}
\end{center}
Most of the ``within-sector'' term is, in this toy, actually the level of a $\kappa$-weighted average
inflation rate moving, not sectors diverging from one another. \textbf{Caveat:} $\bar\pi_t$ here uses
the \emph{welfare} weights $\kappa_i$ (Domar $\times$ rigidity), not the consumption weights
$\bm\beta^H$ that define the true CPI — so $\bar\pi_t$ is a welfare-relevant proxy, not literally
$\pi_t^C$. The two coincide only if $\kappa_i\propto\beta_i^H$, which need not hold since $\kappa_i$
also carries the rigidity ratio $(1-\hat\delta_i)/\hat\delta_i$.

\textbf{Regime decomposition} (toggling how much of the FX shock passes through, as a stand-in for
Float/Managed/Peg):
\begin{center}
\begin{tabular}{lrrrr}
\toprule
Regime & Output-gap & Within-sector & Cross-sector & Total \\
\midrule
Float (full pass-through) & 0.033 & 1.580 & 0.0038 & 1.616 \\
Managed (50\% offset) & 0.027 & 1.530 & 0.0036 & 1.560 \\
Peg (0\% pass-through) & 0.025 & 1.520 & 0.0036 & 1.545 \\
\bottomrule
\end{tabular}
\end{center}

\begin{insight}
\textbf{Why this toy says ``Peg is best,'' while the real project says ``Peg is worst.''} In this
toy, killing FX pass-through only ever \emph{removes} a cost channel — there's no offsetting cost to
pegging. The real model's Peg is dominated instead because pegging doesn't eliminate the
risk-premium/UIP shock, it just closes off the one instrument (the exchange rate) that could
otherwise absorb it, so that shock's variance gets pushed into $\pi_i$ and $y$ instead. This toy
doesn't have a risk-premium shock at all, which is exactly why it can't reproduce that result — a
useful illustration that the shock decomposition step is not optional set-dressing, it's what
determines the regime ranking. See \S\ref{sec:regime}.
\end{insight}

%==============================================================
\section{Scaling to the Real Chile Model}
\label{sec:real}
%==============================================================

The real pipeline is structurally identical to the toy above; only the loadings $L_{i,k}$ come from
solving the full linearized DSGE in Dynare (via the QZ decomposition, \texttt{pol\_fct\_m.m}-style)
instead of being assigned by hand, and there is a genuine third sector plus a fourth shock type
(monetary policy).

\subsection{Where each term is computed}
\begin{center}
\begin{tabular}{lll}
\toprule
Term & File & Object \\
\midrule
Output-gap + within-sector & \texttt{welf.m} / Dynare \texttt{oo\_.var} & \texttt{50*((gamma+phi)*y\^{}2 + pi'*D1*pi)} \\
Cross-sector & \texttt{code/cross\_sector\_welfare.py} & $\Phi_C + \sum_s\lambda_s\Phi_s$, closed form (Remark, \S6) \\
Shock decomposition & \texttt{results/variance\_decomposition.csv} & Dynare's native \texttt{oo\_.variance\_decomposition} \\
\bottomrule
\end{tabular}
\end{center}
\texttt{cross\_sector\_welfare.py} reads the full $3\times3$ unconditional covariance matrix of the
markup gap $\bm\mu_t$ (\texttt{results/markupgap\_chile\_covar.csv}, from reporting variables
\texttt{MARKUPGAP1/2/3} added to the \texttt{.mod} files) rather than just each sector's own
variance — necessary because $\Phi_C,\Phi_s$ are quadratic forms in the \emph{full} covariance
matrix, exactly as in the toy's $\mathrm{Cov}(\mu_1,\mu_2)$ step.

\subsection{Headline numbers (term decomposition, real Chile calibration)}
\begin{center}
\begin{tabular}{lrrr}
\toprule
($\times10^{-4}$) & Float & Managed & Peg \\
\midrule
Output-gap + within-sector & 25.47 & 10.17 & 102.05 \\
Cross-sector (added 2026-07-23) & 0.84 & 0.86 & 4.57 \\
\midrule
\textbf{Total} & \textbf{26.31} & \textbf{11.03} & \textbf{106.62} \\
\bottomrule
\end{tabular}
\end{center}
Same qualitative pattern as the toy: cross-sector is small relative to within-sector (3--8\% here vs.
0.2\% in the toy — the real model's outer CES nest, $\eta=1.5\neq1$, makes it somewhat larger than
the toy's all-Cobb-Douglas case, but still second-order).

\subsection{Shock decomposition (real model)}
\label{sec:regime}
\texttt{results/variance\_decomposition.csv} is Dynare's own forecast-error variance decomposition,
computed the same additive way as \S2's Step 4 (independence of shocks). Reading it for
$\mathtt{y\_gap}$: under Float, the risk-premium shock \texttt{eps\_rp} explains 68.9\% of output-gap
variance; under Peg, 85.8\%; under Managed, 22.2\%. This is the real-model analogue of the toy's
Step 4 table — and it is precisely because \texttt{eps\_rp} dominates so heavily under Peg (not
eliminated by pegging, since UIP/risk-premium is a financial-account object, not an exchange-rate
object) that Peg ends up \emph{worst}, the opposite of the toy's naive "less FX exposure is always
better" intuition. Confirmed directly in the risk-premium-volatility sweep
(\texttt{code/sweep\_risk\_premium\_chile.m}): switching \texttt{eps\_rp} off entirely makes Peg
\emph{roughly tied with Float}, fractionally better — see \texttt{results\_summary.md}.

\subsection{Sector decomposition (real model)}
Mirrors the toy's Step 3 exactly: Domar weight $\lambda_{D,i}$ (size) times a rigidity multiplier
$(1-\hat\delta_i)/\hat\delta_i$ (stickiness) gives $\kappa_i$; combined with each sector's simulated
$\mathrm{Var}(\pi_i)$, Services (the most Domar-weighted, stickiest sector in the Chile calibration)
absorbs the largest share of loss in every regime — the ``Why Services?'' decomposition in
\texttt{code/services\_mechanism\_decomposition.py}, full writeup in \texttt{results\_summary.md}.

\begin{keyresult}
The four decompositions — term, sector, shock, regime — are not four different techniques. They are
the \emph{same} additive-variance identity (Eq.~\ref{eq:linear}) sliced along different axes: sum
over shocks $k$ (term/sector decomposition, holding the object fixed) or sum over objects
(shock decomposition, holding the shock fixed), or recompute the whole object with one loading
zeroed out (regime decomposition). Once you have the loadings $L_{i,k}$ and shock variances
$\sigma_k^2$, every one of these tables is a different reduction of the same matrix
$L_{i,k}^2\sigma_k^2$ — nothing more is needed.
\end{keyresult}

\end{document}
